\documentclass[11pt,a4paper]{article}
\usepackage{fontspec}
\setmainfont{Noto Serif}
\setsansfont{Noto Sans}
\usepackage{babel}
\babelprovide[main,import]{russian}
\usepackage{amsmath,amssymb,mathtools,booktabs,geometry,fancyhdr,microtype,xcolor,hyperref}
\geometry{left=21mm,right=21mm,top=21mm,bottom=21mm}
\definecolor{ink}{HTML}{173D57}
\hypersetup{colorlinks=true,urlcolor=ink,linkcolor=ink,pdftitle={BGS-F. Расстояния между собственными фазами}}
\pagestyle{fancy}\fancyhf{}\lhead{\small BGS-F}\rhead{\small Расстояния между собственными фазами}\cfoot{\thepage}
\setlength{\headheight}{15pt}\setlength{\parindent}{1em}\setlength{\parskip}{2.5pt}
\setlength{\emergencystretch}{3em}\tolerance=3000
\newcommand{\Tr}{\operatorname{Tr}}
\newcommand{\EE}{\mathbb E}
\newcommand{\R}{\mathbb R}
\newcommand{\Z}{\mathbb Z}
\newcommand{\C}{\mathbb C}
\newcommand{\si}{\mathrm{Sine}_2}
\newcommand{\eps}{\varepsilon}
\newcommand{\sinc}{\operatorname{sinc}}
\newcommand{\sect}[1]{\section*{\sffamily\color{ink}#1}}
\newcommand{\subt}[1]{\par\medskip\noindent\textbf{#1.}\quad}
\newcommand{\T}{\mathbb T}
\newcommand{\dist}{\operatorname{dist}}
\begin{document}
\noindent{\small\sffamily\color{ink} МАТЕМАТИЧЕСКАЯ ФИЗИКА\hfill 19 НОЯБРЯ 2026}
\begin{center}{\LARGE\sffamily Расстояния между собственными фазами}\\[4pt]
{\large Конечные отклонения и заполнение арифметических щелей в BGS-F}
\end{center}
Пусть $e^{2\pi i\theta_j}$ --- собственные числа исходной матрицы $U_N=e^{i\tau NA_N}M_N$, $\tau=10^{-80}$. Среднее расстояние между фазами на окружности длины 1 равно $1/N$. Изучается та же последовательность по всем нечётным $N$, с единственным случайным параметром --- общим сдвигом начала отсчёта. Точная матрица и символ выписаны в начале следующей части и в постановке задачи внутри архива.

\sect{Результаты}
\textbf{Все нечётные $3\le N\le5\cdot10^{39}$.} Для любой пары фаз развёрнутое расстояние $N\|\theta_i-\theta_j\|_{\T}$ лежит вне $(0.10,0.65)$. Построен один ограниченный непрерывный функционал точечного процесса, который равен нулю при каждой реализации общего сдвига во всём этом диапазоне. Его ожидание для $\si$ превышает $0.0026$.

\textbf{Полный парный тест.} Для каждого из 1034 нечётных кратных пеллева числа $h_{175}$ на отрезке $[10^{70},2\cdot10^{70}]$ получено строгое неравенство
\[
                        \boxed{\Delta_{2,N}(b_*)>221.}
\]
Здесь участвует исходный вес и вся сумма по $1\le m<2N$. Для этих же размерностей вероятность обнаружить хотя бы один уровень в случайном окне длиной в одну среднюю щель меньше $0.000073$. При $N=h_{207}>8.58\cdot10^{78}$ ожидание ограниченного непрерывного теста окна отличается от синусного эталона более чем на $.0465$. Это конечные утверждения.

\textbf{Бесконечная последовательность исходных матриц.} При $N=h_k$, $k\to\infty$ нечётно, каждая щель повёрнутой сетки из $k$ равноотстоящих точек получает не менее $cN/k$ фаз на расстоянии $\ge c/k$ от сетки с некоторой постоянной $c>0$. Поворот может зависеть от $N$. Кроме того, доля фаз на расстоянии $\ge1/(6k)$ не меньше $1/3-o(1)$. Оптимальная максимальная длина переноса от невозмущённого спектра к возмущённому имеет порядок $1/k$. Здесь используются логарифмическая следовая оценка и строгое неравенство давления предыдущей части; численная граница их асимптотической применимости неизвестна.

\textbf{Граница вывода.} Ни конечное нарушение (P), ни заполнение щелей шириной порядка $1/\log N$ не определяют предельный закон при ширине $1/N$. Ниже также получены оценки скорости собственных фаз и явный асимптотический закон первого следа. Весь сводный документ включает предыдущие три отчёта с их собственной нумерацией.
\newpage
\sect{1. Запрещённый интервал и тест слабой сходимости}
Положим $o_N=\operatorname{ord}_N(S)$, где $S=\left(\begin{smallmatrix}3&4\\2&3\end{smallmatrix}\right)$. По ковариантности Вейля $M_N^{o_N}$ скалярна. Поэтому все её исходные фазы лежат в некотором повороте решётки $o_N^{-1}\Z/\Z$. Из расщепления $a+b\sqrt2$ при простых модулях и подъёма порядка (раздел 16 следующей части) имеем $o_N\le4N/3$.

Рассмотрим путь $U(t)=e^{itA}M$, $0\le t\le\xi=\tau N$. У его собственных фаз есть непрерывные вещественные поднятия $\theta_j(t)$, каждое из которых абсолютно непрерывно и удовлетворяет $2\pi|\dot\theta_j|\le\|A\|$. Это следует локально из аналитического эрмитова логарифма унитарного пути с разрезом вне спектра и теоремы Реллиха. Дифференцирование в регулярной ветви даёт $2\pi\dot\theta_j=\langle u_j,Au_j\rangle$; в кратном кластере скорости задаются сжатием $A$. Локальные ветви продолжаются по компактному интервалу. Из 32 коэффициентов исходного символа $\|A\|\le C=593/500$. Значит, все фазы остаются в окрестности указанной сетки радиуса $C\tau N/(2\pi)$, с одновременным сопоставлением по кратностям.

Разность двух фаз отличается от элемента $o_N^{-1}\Z/\Z$ не более чем на $C\tau N/\pi$. При $N\le5\cdot10^{39}$ после умножения на $N$ это строго меньше $0.10$. Ненулевое расстояние в решётке после развёртки не меньше $N/o_N\ge0.75$. Следовательно, любая пара либо имеет расстояние меньше $0.10$, либо больше $0.65$.

Пусть $f$ равно 1 на $[-.05,.05]$, линейно убывает до нуля к $\pm.06$ и равно нулю дальше. Положим $g(x)=f(x-.5)$,
\[
 F(\zeta)=\min(1,\langle f,\zeta\rangle)\min(1,\langle g,\zeta\rangle).
\]
Функционал ограничен и непрерывен в слабой топологии локально конечных мер при тестировании непрерывными финитными функциями. Любые две точки из носителей имеют расстояние в $[.38,.62]$. Таким образом $F(\Xi_{N,\alpha})=0$ при всех $\alpha$. Для $\si$ обозначим через $X,Y$ числа точек в двух интервалах, где $f,g$ равны единице. Для целых $X,Y\ge0$
\[
 F\ge\mathbf1_{X,Y\ge1}\ge XY-\binom X2Y-X\binom Y2 .
\]
Второе неравенство следует подстановкой целых значений (правая часть $XY(4-X-Y)/2$). На этих интервалах $|x-y|\ge.4$, поэтому $\rho_2(x,y)\ge1-(2.5/\pi)^2$. По неравенству Адамара $\rho_3\le1$, поскольку ядро синуса положительно полуопределено с единичной диагональю. Интегрируя, получаем
\[
                 \EE_{\si} F\ge.01\{1-(2.5/\pi)^2\}-.001>.0026 .
\]
Так исключено совпадение с эталоном для ограниченной непрерывной статистики во всём указанном конечном диапазоне.
\newpage
\sect{2. От покрытия спектра к полному тесту (P)}
Пусть $b=b_*$, $w=|\widehat b|^2$, $\beta=\int b$. Исходный тест записывается как
\[
 \Delta_{2,N}=N^{-2}\sum_{m=1}^{2N-1}w(m/N)(|p_{m,N}|^2-\min(m,N)),
 \qquad p_{m,N}=\Tr U_N^m.
\]
Если $B(\alpha)=\sum_{j,l}b(N(\theta_j-\alpha+l))$, периодизация и Парсеваль дают $\EE B=\beta$ и $\mathrm{Var}\, B=2N^{-2}\sum_{m\ge1}w(m/N)|p_{m,N}|^2$. Возьмём объединение дуг $E$ вокруг $o_N$ узлов радиуса $C\tau N/(2\pi)+1/(4N)$. Каждая дуга радиуса $1/(4N)$ вокруг каждой собственной фазы включена в $E$. Положительность $b$ даёт
\[
 |E|\le Q_N:=\frac{o_N}{N}\left(\frac12+\frac{C\tau N^2}{\pi}\right),
 \qquad \int_E B(\alpha)d\alpha\ge\int_{-1/4}^{1/4}b(x)dx>0.18 .
\]
Действительно, нормированный вес $\phi(t)\,dt/I$ симметричен и убывает по $|t|$, поэтому его второй момент не выше $1/3$ (распределение $|t|$ стохастически меньше равномерного на $[0,1]$). Из $\cos y\ge1-y^2/2$ на $|x|\le1/4$ имеем $\check\phi(x)/I\ge1-2\pi^2x^2/3>0$. Интегрирование $b(x)\ge(1-4\pi^2x^2/3)/2$ даёт $\frac14(1-\pi^2/36)>.18$. Также $\beta=\int\phi^2/(2I^2)\le e^{-1}/(2I)<.7$, поскольку $I\ge e^{-4/3}$, а $|\widehat b|\le\beta$. Неравенство Коши и ограничение нормированного вычитаемого члена сверху числом $2\beta^2$ дают
\[
 \boxed{\displaystyle\Delta_{2,N}\ge\frac{.18^2}{2Q_N}-\tfrac52(.7)^2.} \tag{1}
\]
Точно так же вероятность попадания случайного окна длины $1/N$ в спектр ограничена $Q_N+o_N/(2N)$, по покрытию с добавкой $1/(2N)$.

Пусть $D\mid N$, $S^k=I\bmod D$. Тогда $o_N/N\le4k/(3D)$. В самом деле, усилим простые степени $D$ до всех их степеней в $N$, получив $D'$. По биномиальной формуле при каждом $p^a\parallel D$ дополнительный множитель порядка не превышает $p^{v_p(D')-a}$. Следовательно $o_{D'}\mid kD'/D$. Сопростой остаток $n=N/D'$ имеет порядок не более $4n/3$ (для $n=1$ берём порядок 1). Китайская теорема завершает оценку. Для $D=h_k$ с нечётным $k$ условие на $S^k$ следует из $(1+\sqrt2)^{2k}$; оно дополнительно проверено целыми числами для всей таблицы.
\begin{center}\small\begin{tabular}{@{}rrrrr@{}}\toprule
$k$ & Интервал для $N$ & Число нечётных кратных $h_k$ & $\Delta_{2,N}>$ & Вероятность попадания $<$\\\midrule
13 & $[3,10^{39}]$ & $>10^{34}$ & 86 & .000368\\
175 & $[10^{70},2\cdot10^{70}]$ & 1034 & 221 & .000073\\
185 & $[10^{72},2\cdot10^{72}]$ & 16 & 140 & .000115\\
191 & $[10^{73},2\cdot10^{73}]$ & 1 & 270 & .000060\\\bottomrule\end{tabular}\end{center}
Подсчёт делается целочисленным делением концов интервала. В частности, за второй строкой стоит $N$ порядка $10^{70}$, при котором $\tau N^2$ порядка $10^{60}$: малость перемещения в единицах средней щели здесь не используется.
\newpage
\sect{3. Скорость фаз и возврат при $N=h_{207}$}
Обозначим $W_t=\operatorname{Ad}U(t)$, $W_0=\operatorname{Ad}M$, а $\|X\|_n^2=N^{-1}\Tr X^*X$. Для любого $L$ и $T=2L+1$ в собственном базисе $U(t)$ диагонали $A$ и $T^{-1}\sum_{j=-L}^LW_t^jA$ совпадают. Поэтому
\[
 N^{-1}\sum_s|2\pi\dot\theta_s(t)|^2
 \le\left\|T^{-1}\sum_{j=-L}^LW_t^jA\right\|_n^2 .                       \tag{2}
\]
Выбираем $L=94$. Для $t\in[0,\tau N]$ точный коммутатор с модой Вейля даёт
$\|(\operatorname{Ad}e^{itA}-I)T(v)\|_n\le80\tau n(v)$, $n(a,b)=|a|+\sqrt2|b|$, как в разделе 1 следующей части. Справедлива та же оценка с минусом перед $t$. В телескопической разности $W_t^j A-W_0^jA$ сравниваем на чистых модах $S^i d$ или $S^{-i}d$. Для них $n\le14\Lambda^i$, $\Lambda=3+2\sqrt2$, а сумма модулей коэффициентов $\le C$. Вплоть до $|j|=94$ разность в норме не более
\[
 e_0=80\tau C\,14\,\frac{\Lambda(\Lambda^{94}-1)}{\Lambda-1}<0.00001469 .
\]
Среднее от $-94$ до 94 при $t=0$ можно сдвинуть по времени к $0,\ldots,188$, сохраняя норму. Складываем целые гауссовы коэффициенты $4000c_d$ в ячейках $S^id\bmod N$: сумма квадратов модулей, делённая на $(4000T)^2$, точно равна квадрату нормы среднего. Для $k=197,203,205,207,209$, $N=h_k$ значение одинаково: $44571523/20412000000$. Отсюда правая часть (2) не превосходит $v^2$ при $v=0.046744$. Интегрирование по ветвям (неравенство Минковского) даёт среднеквадратичное отклонение в радианах не более $\xi v$.

При $M^k=e^{i\gamma}I$ имеем
$N^{-1}\Re(e^{-i\gamma}p_{k,N})\ge1-(k\xi v)^2/2$ из $\cos x\ge1-x^2/2$. Для $k=207$ с неокруглённой верхней границей $v$ это больше $0.65530$. Достаточно одного положительного слагаемого с $m=k$, чтобы получить
\[
             \Delta_{2,h_{207}}(b_*)>0.02448 .
\]
Здесь $w(k/N)>.10892$, а полный вычитаемый член меньше $.022292$. Для проверки $I$ и свёртки скрипт делит $[-1,1]$ на 256 ячеек и вычисляет суммы Дарбу. На узлах используются $(1-x/2048)^{2048}\le e^{-x}\le(1+x/2048)^{-2048}$, затем округление наружу на сетку $2^{-40}$. Свёртка чётных убывающих весов убывает при $t\ge0$ по представлению через интервалы. Сгруппировав частоты в ячейки ширины $1/128$, берём левый конец для $w$, правый для рампы и не более $N/128+1$ частот на ячейку. Аналогично получены $\Delta_{2,h_{203}}>.0865$ и $\Delta_{2,h_{205}}>.0844$. При $h_{209}$ оценка скорости перестаёт давать положительную нижнюю границу нормированного вещественного следа. Старый сертификат в конце сводного PDF получен другим методом.
\newpage
\sect{3а. Локальная занятость при $N=h_{207}$}
В той же размерности из (2) следует отклонение для \emph{ограниченного непрерывного} наблюдения почти единичного окна. Пусть $f_1$ равно 1 на $[-1/2,1/2]$, равно нулю вне $[-.505,.505]$ и линейно на двух переходных отрезках. Тогда
\[
 F_1(\zeta)=\min(1,\langle f_1,\zeta\rangle),\qquad
 \boxed{\ \EE_{\si}F_1-\EE_\alpha F_1(\Xi_{h_{207},\alpha})>.0465 .\ }     \tag{2а}
\]
Здесь $h_{207}>8.58\cdot10^{78}$, и грубая граница из раздела 1 не контролирует попадания в это окно.

Действительно, оценка скорости даёт $D=k\xi v/(2\pi)<.13215$ для среднеквадратичного сдвига в единицах периода решётки $1/k$. Доля фаз на расстоянии больше $a/k$ от ближайшего исходного узла не превосходит $(D/a)^2$. Если окно длины $l/N$ встречает хотя бы одну из остальных фаз, его центр лежит в объединении $k$ интервалов общей длины не более $2a+kl/N$. Для исключительных фаз вероятность попадания не более $l(D/a)^2$, по сумме вероятностей попадания в каждую из них. Следовательно при $a=.26$, $l=1.01$
\[
 \EE_\alpha F_1(\Xi_{h_{207},\alpha})\le 2a+kl/N+l(D/a)^2<.78091.
\]

Для $\si$ пусть $X$ --- число точек в центральном интервале длиной 1. Тогда $F_1\ge\mathbf1_{X\ge1}\ge X-\binom X2$. По парной функции синус-ядра
\[
 \EE_{\si}F_1\ge\tfrac12+\int_0^1(1-s)
       \left(\frac{\sin\pi s}{\pi s}\right)^2ds>.82749.
\]
Нижняя граница интеграла получена без численного интегрирования. В $(1-\cos 2x)/(2x^2)$ заменяем косинус сверху его многочленом Тейлора степени 12; неравенство следует последовательным интегрированием $1-\cos u\ge0$ с нулевыми начальными остатками. После интегрирования по $s$ шесть членов равны
\[
 \sum_{j=1}^6\frac{(-1)^{j+1}\,2^{2j-1}\pi^{2j-2}}{(2j)!(2j-1)(2j)} .
\]
Скрипт подставляет нижнюю или верхнюю рациональную границу $\pi$ по знаку слагаемого и вычитает верхнюю границу, полученную для квантовой карты. Таким образом, конечный порог (2а) не опирается на переход от неограниченного второго момента к слабой сходимости.
\newpage
\sect{4. Короткие интервалы при всех больших $N$}
Пусть $\mu_N=N^{-1}\sum_j\delta_{\theta_j}$, $h=\log\Lambda$, $p=P_{\rm top}(F,-\frac12\log J^uF)$, так что $0<p<h/2$. Используем доказанную в части о давлении оценку: для каждого $a<1/p$ при $H=\lfloor a\log N\rfloor$ существует $\eta>0$ такое, что $\sum_{1\le m\le H}|p_{m,N}/N|^2=O(N^{-2\eta})$. Достаточно взять $\epsilon>0$ с $a(p+\epsilon)<1$ в доказательстве той части. Там же указаны используемые логарифмическая формула Фора [1] и адаптация к нечётному $N$.

Для любой дуги $I$ на окружности периодические многочлены Сельберга--Ваалера [2] степени $H$ ограничивают её индикатор сверху и снизу и имеют интегралы $|I|\pm1/(H+1)$. Можно пользоваться открытыми, замкнутыми и полуоткрытыми концами: в граничной точке мажоранта достигает не меньше 1, миноранта не превосходит 0. Эти свойства получаются из аппроксимации периодической функции $x-\lfloor x\rfloor-1/2$ в двух концах с ошибкой, ограниченной сдвигами ядра Фейера, делёнными на $2(H+1)$. У каждой мажоранты $L^1$-норма ограничена 2; у миноранты --- 4 по разности мажоранты и миноранты. Все коэффициенты поэтому равномерно ограничены. Интегрирование по $\mu_N$, неравенство Коши и затем $a\uparrow1/p$ дают
\[
 \limsup_{N\to\infty}(\log N)\sup_I|\mu_N(I)-|I||\le p.                  \tag{3}
\]
В частности, на \emph{любом} интервале разностей фаз длины $s/N$, с фиксированным $s>0$, число упорядоченных пар, включая диагональ, не превосходит $(p+o(1))N^2/\log N$. Фиксируем первую фазу и применяем (3) ко второй. Это контроль приблизительных совпадений, равномерный по положению окна. Для статистики с единичной интенсивностью в микроскопическом масштабе требуется точность порядка $N$.

Оценка коротких следов ограничивает также скопления сильных возвратов при сколь угодно больших временах. В пространстве $L^2(\mu_N)$ возьмём функции $e^{2\pi ijx}$, $0\le j\le H$. Их матрица Грама отличается от единичной по операторной норме не более чем на $2\sum_{1\le j\le H}|p_{j,N}|/N=o(1)$ (оценка строк). Применяя сопряжённый оператор анализа к единичному вектору $e^{-2\pi itx}$, для каждого целого $t$ имеем
\[
           \sum_{j=0}^{H}|p_{t+j,N}|^2/N^2\le1+o(1)                     \tag{4}
\]
с одной и той же ошибкой для всех $t$. Для любого $b>0$ в полосе $N\le m\le2N$ доля времён с $|p_{m,N}|\ge bN$ сверху не более $(p+o(1))/(b^2\log N)$, по разбиению на блоки. Возможность отдельного сильного возврата при $m$ порядка $N$ сохраняется.
\newpage
\sect{5. Каждая щель пеллевой решётки заполняется}
Пусть $N=h_k$, $k\to\infty$ по нечётным числам; тогда $k\sim2\log N/h$, $M^k$ скалярна. Для произвольного поворота $G_\gamma=\gamma+k^{-1}\Z/\Z$ из (3) следует
\[
 \limsup k\sup_I|\mu_N(I)-|I||\le r:=2p/h<1.
\]
Пусть $c=(1-r)/8>0$. В каждом из $k$ промежутков между соседними узлами возьмём его середину в виде дуги длины $(1-2c)/k$. В ней при всех достаточно больших $k$ лежит не меньше $cN/k$ собственных фаз, и расстояние каждой из них от $G_\gamma$ не меньше $c/k$. Оценка равномерна по всем поворотам и промежуткам. У невозмущённой матрицы все фазы лежали на узлах своей решётки.

Количественно более далёкое перемещение следует из $p_{k,N}/N=o(1)$. Если $q$ --- доля уровней на расстоянии $\ge1/(6k)$ от $G_\gamma$, то
\[
 \Re(e^{-2\pi ik\gamma}p_{k,N}/N)\ge\tfrac12(1-q)-q,
 \qquad q\ge1/3-o(1).
\]
Тот же вывод локально верен внутри любой дуги $I=I_N$ при $k|I|\to\infty$: доля берётся среди находящихся в дуге уровней. Выберем $a$ с $2/h<a<1/p$, $H=\lfloor a\log N\rfloor$ и $B=\lfloor\min((H-k)/2,k/2)\rfloor\asymp k$. Для $I$ возьмём среднее двух экстремальных многочленов степени $B$ из раздела 4. Ошибка аппроксимации в $L^1(\mu_N)$ ограничена интегралом полуразности, то есть $1/(B+1)+O(\sqrt H N^{-\eta})$. После умножения среднего на $e^{2\pi ikx}$ все частоты лежат в $[1,H]$, его интеграл по $\mu_N$ имеет модуль $O(\sqrt H N^{-\eta})$. Получаем $|\int_I e^{2\pi ikx}d\mu_N|=O(1/k)$. По (3) $\mu_N(I)\sim|I|$, что и даёт локальный вывод.

\subt{Цена переноса спектра} Обозначим через $W_\infty$ наименьшее максимальное круговое расстояние при сопоставлении двух списков фаз с кратностями (для общих мер --- минимальная существенная верхняя грань расстояния в сцеплении). Для равновесных эмпирических мер существование сцепления с радиусом $d$ даёт перестановку с тем же радиусом по теореме Холла. Если $D=\sup_I|\mu(I)-|I||$, то $W_\infty(\mu,dx)\le D/2$: в периодическом поднятии функции распределения $A(x)$ колебание $A(x)-x$ равно $D$, учитывая левые и правые пределы. Сдвигаем равномерный параметр на середину этого диапазона, затем пользуемся обобщённой обратной $A^{-1}$; внутри каждого скачка оба его конца удовлетворяют тому же ограничению.

Для $1\le m<k$ чистый след ограничен $|\Tr M^m|\le\Lambda^{\min(m,k-m)/2}=O(\sqrt N)$ по числу неподвижных мод Вейля и определителю $S^m-I$. Обратное преобразование Фурье на $k$ узлах даёт кратность $N/k+O(\sqrt N)$ на каждом узле. Отсюда для любой дуги невязка меры чистого спектра с длиной не более $1/k+O(k/\sqrt N)$. Для его произвольного поворота $\mu_M^\gamma$ неравенство треугольника с равномерной мерой и (3) теперь дают
\[
 \tfrac16\le\liminf k W_\infty(\mu_N,\mu_M^\gamma)
 \le\limsup k W_\infty(\mu_N,\mu_M^\gamma)\le\tfrac12+p/h<1.               \tag{5}
\]
Поворот может зависеть от $k$. Нижняя оценка следует из доли ушедших от решётки. Таким образом, порядок перемещения измерен с обеих сторон на бесконечном ряду исходных размерностей.
\newpage
\sect{6. Устойчивость локальных наблюдений по параметру}
Усилим оценку (2) в асимптотическом режиме. При всех $|t|\le\tau N$ имеется равномерная граница
\[
            N^{-1}\sum_j|2\pi\dot\theta_j(t)|^2=O(1/\log N) .             \tag{6}
\]
Константа в (6) не используется для конечных сертификатов. Приведём обоснование для точного оператора в вейлевом базисе. В разделах 1--2 следующей части для него получены на поднятой решётке матрица $L=E\mathcal S$, $\mathcal S e_v=e_{Sv}$, и оценки
\[
 \sum_w |E_{wz}|e^{2n(w-z)}\le e^{10^{-64}n(z)},\qquad
 \sum_w |L_{wv}|e^{g(w)-g(v)}\le e^{-.8n(v)}\quad(v\ne0),
\]
где $g=|a-\sqrt2 b|-|a+\sqrt2 b|$. Их доказательство использует лишь $|t|/N\le10^{-80}$. Изоляция нулевой моды и отсутствие постоянного члена в $A$ при больших $N$ существенны.

Берём $R_0=100\log N$, $R_j=8^j R_0$, $T=\lfloor\log N/(4\log8)\rfloor$. При $j\le T$ эти радиусы меньше $N/16$ для больших $N$. За шаг из $n(v)\le R_j$ в центрированные классы $n(w)>R_{j+1}$ все поднятия выходной метки имеют $n(w+N\ell-Sv)>2R_j$, поскольку $n(Sv)<6R_j$. Первая из двух оценок даёт норму хвоста $O(N^2e^{-3R_j})$ после периодизации. Телескопированием можно на всех $j<T$ заменить эволюцию исходного $A$ на степени сжатия внутри $N/16$ с ошибкой $O(T N^2e^{-3R_0})$: усечённая после каждого шага по нарастающим радиусам цепь сравнивается отдельно с полной и сжатой эволюцией. В сжатии стоимость замены периодической матрицы на поднятую равна $O(Te^{-N})$ по оценке (5) следующей части.

Вторая оценка сжимает взвешенную $\ell^1$-норму на бесследовых метках в $q=e^{-.8}$ раз. Вход и тестовый выход $A$ имеют фиксированный конечный носитель, на котором вес и его обратный ограничены независимо от $N,t$. Следовательно $|\langle A,W_t^j A\rangle_n|\le Cq^j+O(N^{-20})$ при $1\le j<T$. Для отрицательных $j$ используется сопряжение. Квадрат нормы среднего по $0,\ldots,T-1$ выражается через сумму этих корреляций с треугольными весами и не превосходит $C'/T$. Равенство диагоналей из раздела 3 доказывает (6).

Отсюда при двух силах кика $u,u'$ в $[-\tau,\tau]$ среднеквадратичное расстояние между поднятыми фазами вдоль пути ограничено $C N|u'-u|/\sqrt{\log N}$. Если сгладить локальный процесс любой $f\in C_c^1(\R)$ и использовать один равномерный сдвиг $\alpha$ в обоих процессах, то
\[
 \EE_\alpha|\langle f,\Xi_{u,N}-\Xi_{u',N}\rangle|
 \le C\|f'\|_1 N^2|u'-u|/\sqrt{\log N}.
\]
Действительно, интегрирование по $\alpha$ и суммирование по периодам даёт на одну ветвь оценку $\|f'\|_1|\theta_j-\theta'_j|$. Это указывает ширину изменения параметра, при которой локальное наблюдение устойчиво в среднем, в частности при $u=\tau$ и $|u'-u|=o(\sqrt{\log N}/N^2)$. Во всех результатах для (P) выше конечная сила по-прежнему в точности $\tau$.
\newpage
\sect{7. Явный закон интенсивности первого следа}
Для полноты вычислим один спектральный предел при фиксированной силе непосредственно по четырём неподвижным точкам. Он показывает, что уже первая гармоника имеет собственный арифметический закон. Пусть индексы $a,b=0,1$ нумеруют точки $z_{ab}=(a/2,b/2)$ в координатах $(q,p)$, и
\[
 (c_{00},c_{01},c_{10},c_{11})=\tfrac1{125}(-19606,-19594,24306,24894),
 \quad d_{ab}=(4-2\pi\tau c_{ab})^{-1/2}.
\]
Обозначим $A_*=d_{00}-d_{01}$, $B_*=d_{10}+d_{11}$. Тогда по всем нечётным $N$
\[
 \frac{\Tr U_N}{\Tr M_N}
 =i^N A_*e^{9i\tau N/10}+B_* e^{-11i\tau N/10}+O_\tau(N^{-1}).           \tag{7}
\]
Здесь $|\Tr M_N|=1$ по квадратичной сумме Гаусса с обратимым старшим коэффициентом. Распределение $|\Tr U_N|^2$ при эмпирическом подсчёте нечётных $N\le X$, $X\to\infty$, имеет предел с арксинусной плотностью на $[(B_*-|A_*|)^2,(B_*+|A_*|)^2]$. Действительно, в каждом нечётном классе по модулю 4 фаза $2\tau N$ равномерно распределена по модулю $2\pi$ (геометрические прогрессии по критерию Вейля, поскольку $\tau/\pi$ иррационально). Применяем (7); $A_*\ne0$. Этот предел по размерам характеризует только указанную последовательность чисел.

Проверим фазы в (7). Квадратичное дополнение в точной модели даёт для $v=(a,b)\in\Z^2$
\[
 \Tr(T(v)M)/\Tr M=e_N(\bar2 b^2-\bar4 a^2).
\]
Для физической моды $(-1)^{ab}T(v)$ это равно
$(-1)^{ab+b} i^{Na^2}\exp(2\pi i(b^2/2-a^2/4)/N)$.
Преобразование Фурье первого множителя по двум чётностям даёт для локализаций соответственно $i^N/2,-i^N/2,1/2,1/2$. Для фиксированного гладкого символа это вычисляет главный коэффициент локализованного следа при $\tau=0$: разворачиваем быстро сходящийся ряд Фурье и оцениваем вклад последней экспоненты с ошибкой $O(N^{-1})$.

Все четыре точки остаются неподвижными при включении кика: $\nabla g=0$; топологическое сопряжение с линейной картой исключает дополнительные неподвижные точки. В них $g=9/10$ при $q=0$ и $-11/10$ при $q=1/2$, $g_{qq}/\pi^2=c_{ab}$, $g_{qp}=g_{pp}=0$. Поэтому $|\det(I-DF)|=4-2\pi\tau c_{ab}$. Применим одношаговую локализованную следовую формулу [1]: главный модуль есть обратный квадратный корень определителя. Вдоль пути включения кика определитель не обращается в нуль, так что индекс Маслова сохраняется. Вейлев генератор не имеет субглавного члена, а классический путь кика в этих точках постоянен; приращение экспоненты действия равно $iN\tau g(z_{ab})$. Начальные локализованные знаки уже вычислены. Это даёт (7) без предположения о случайных фазах.

В частности, $(B_*-|A_*|)^2>1+300\tau$. По выпуклости $B_*\ge1+\pi(c_{10}+c_{11})\tau/8>1+154\tau$, тогда как $|A_*|\le\pi|c_{00}-c_{01}|\tau/8<.04\tau$. Следовательно, для всех достаточно больших нечётных $N$ имеем $|\Tr U_N|^2>1+300\tau$. В полном (P) вес этого первого времени имеет порядок $N^{-2}$.
\newpage
\sect{8. Предел точности и воспроизведение}
Для исходного процесса установлены различные масштабы. На конечном диапазоне из раздела 1 исключены пары в целом интервале размеров, измеряемых в средних щелях. В разделах 2 и 3 доказаны строгие отклонения исходного (P) без каких-либо замен его веса. При $N=h_k\to\infty$ каждая прежняя арифметическая щель получает уровни внутри; характерный масштаб такого сдвига составляет $N/k$ средних щелей. Подсчёт пар на расстоянии $s/N$ при всех больших $N$ в разделе 4 ограничивает их число сверху величиной порядка $N^2/\log N$. Для получения парной корреляции нужно вычислить этот подсчёт после сглаживания на масштабе порядка $N$. Утверждения о числе пар и о перемещении уровней допускают разные микроскопические законы. Основная гипотеза BGS-F и предел (P) остаются нерешёнными.

\subt{Целочисленные сертификаты} Файл \texttt{spacings/certify.py} проверяет порог запрета, нижнюю границу ожидания $F$ для $\si$, полные границы для четырёх прогрессий и пять норм орбитных средних. Все неравенства выполняются с рациональными числами. Концы направленных округлений сохранены в JSON. Для $\pi$ использована формула Мэчина с чередующимися остатками. Скрипт \texttt{harmonic.py} точно вычисляет значения символа, гессианы и четыре весовых коэффициента Фурье.

\subt{Независимая матричная проверка} Скрипт \texttt{diagnostics.py} строит маленькие матрицы непосредственно по трём диагоналям кика и ядру $M$. Проверяется соответствие фаз на окружности, равенство двух вычислений нормы орбитного среднего и коэффициент 2 в формуле Парсеваля. Здесь использовано увеличенное тестовое $t=.24/N$ для видимого смещения, только в численной проверке тождеств. В разделе 7 для дополнительной диагностики можно сравнить (7) с прямым вычислением следа при тестовой силе $10^{-4}$ (\texttt{harmonic\_diagnostic.py}); эта сила также не входит в результаты для исходной карты.

\subt{Аналитические зависимости} Оценки (3)--(5) используют доказательство строгого неравенства давления и применения полуклассической формулы из предшествующего отчёта. Константа $c=(1-2p/h)/8$ положительна; гарантированная численная нижняя граница для неё при $10^{-80}$ здесь не вычислена. Никакая из этих асимптотических констант не подставляется при проверке конечных порогов. Доказательства (1) и (2) и вычислительные сертификаты логарифмической следовой формулы не требуют.

\subt{Источники} [1] F. Faure, \emph{Semi-classical formula beyond the Ehrenfest time in quantum chaos. (I) Trace formula}, Ann. Inst. Fourier 57 (2007), 2525--2599, \href{https://doi.org/10.5802/aif.2341}{DOI 10.5802/aif.2341}. [2] J. D. Vaaler, \emph{Some extremal functions in Fourier analysis}, Bull. Amer. Math. Soc. 12 (1985), 183--216, \href{https://doi.org/10.1090/S0273-0979-1985-15349-2}{DOI 10.1090/S0273-0979-1985-15349-2}. Исходная постановка, три предыдущих отчёта и их расчёты включены в сводный пакет.
\end{document}
